% ============================================================
% 快速验证：编译所有 8 个 TikZ 流程图场景到独立的 PDF
% 用法：xelatex 本文件，自动生成 8 页 PDF（每页一个场景）
%       也可取某一页做单页 PDF：如 gs -sDEVICE=pdfwrite -dFirstPage=3 -dLastPage=3 本文件
% ============================================================
\documentclass[a4paper]{article}
\usepackage[margin=0.5cm,paperwidth=21cm,paperheight=29.7cm]{geometry}
\pagestyle{empty}
\usepackage{fontspec}
\setmainfont{Times New Roman}
\setsansfont{Arial}
\usepackage{ctex}
\newcommand{\fontdir}{../论文写作/LaTeX模板_2026国赛版/fonts/}
\setCJKfamilyfont{song}[%
  Path = \fontdir,
  UprightFont = SourceHanSerifCN-Regular.otf,
  BoldFont   = SourceHanSerifCN-Bold.otf,
  AutoFakeBold = true,
]{SourceHanSerifCN}
\newcommand*{\song}{\CJKfamily{song}}
\usepackage{tikz}
\usetikzlibrary{arrows.meta,positioning,calc,shapes.geometric,shapes.multipart,backgrounds,fit,decorations.pathreplacing}

% ---- 颜色方案 ----
\definecolor{primary}{HTML}{2B579A}
\definecolor{primaryDark}{HTML}{1A3A6E}
\definecolor{secondary}{HTML}{D9E4F2}
\definecolor{accent}{HTML}{D35400}
\definecolor{accentLight}{HTML}{FDEBD0}
\definecolor{decisionFill}{HTML}{E8F8F5}
\definecolor{decisionBorder}{HTML}{117A65}
\definecolor{gray}{HTML}{95A5A6}
\definecolor{textDark}{HTML}{2C3E50}
\definecolor{arrowColor}{HTML}{5D6D7E}
\definecolor{groupBg}{HTML}{F4F6F7}

% ---- 节点样式 ----
\tikzset{
  process/.style={rectangle,draw=primaryDark,line width=0.8pt,fill=primary,text=white,minimum width=3.8cm,minimum height=1.0cm,align=center,font=\song\footnotesize\bfseries,inner sep=6pt,rounded corners=2pt},
  subprocess/.style={rectangle,draw=primaryDark,line width=0.7pt,fill=secondary,text=textDark,minimum width=3.2cm,minimum height=0.85cm,align=center,font=\song\footnotesize,inner sep=4pt,rounded corners=2pt},
  startstop/.style={rectangle,draw=accent,line width=1.0pt,fill=accentLight,text=accent!85!black,minimum width=3.2cm,minimum height=0.9cm,align=center,font=\song\footnotesize\bfseries,inner sep=6pt,rounded corners=8pt},
  decision/.style={diamond,draw=decisionBorder,line width=0.9pt,fill=decisionFill,text=textDark,minimum width=3.2cm,minimum height=1.2cm,align=center,font=\song\footnotesize\bfseries,inner sep=2pt,aspect=2.0},
  data/.style={trapezium,trapezium left angle=70,trapezium right angle=70,draw=primaryDark,line width=0.7pt,fill=secondary,text=textDark,minimum width=3.2cm,minimum height=0.9cm,align=center,font=\song\footnotesize,inner sep=6pt},
  annotation/.style={rectangle,draw=none,fill=none,text=gray,font=\song\scriptsize\itshape,align=left},
  groupbox/.style={rectangle,draw=gray,line width=0.6pt,dashed,fill=groupBg,fill opacity=0.5,inner sep=10pt,rounded corners=6pt},
}
\tikzset{
  arrowMain/.style={-{Stealth[length=2.2mm,width=2.0mm]},line width=0.9pt,draw=arrowColor},
  arrowSub/.style={-{Stealth[length=1.8mm,width=1.6mm]},line width=0.7pt,draw=gray},
  arrowBidir/.style={{Stealth[length=2.0mm,width=1.8mm]}-{Stealth[length=2.0mm,width=1.8mm]},line width=0.8pt,draw=arrowColor},
  arrowDashed/.style={-{Stealth[length=2.0mm,width=1.8mm]},line width=0.7pt,draw=gray,densely dashed},
  lineNoArrow/.style={line width=0.7pt,draw=arrowColor},
}
\tikzset{
  edgeLabel/.style={font=\song\footnotesize,fill=white,fill opacity=0.85,inner sep=2pt,text=textDark},
  edgeLabelYes/.style={edgeLabel,text=decisionBorder},
  edgeLabelNo/.style={edgeLabel,text=accent},
}
\newcommand{\nodeDist}{1.2cm}
\newcommand{\nodeDistWide}{1.8cm}
\newcommand{\colDist}{1.5cm}

\newcommand{\sceneTitle}[1]{{\large\bfseries\song #1}\vspace{0.5cm}}

\begin{document}

% ---- 场景 1：顺序流程 ----
\sceneTitle{场景 1：顺序流程图}
\begin{center}
\begin{tikzpicture}[node distance=\nodeDist]
  \node[startstop] (N1) {数据预处理与清洗};
  \node[process, below=of N1] (N2) {特征提取与选择};
  \node[process, below=of N2] (N3) {模型建立与参数估计};
  \node[process, below=of N3] (N4) {模型求解与优化};
  \node[startstop, below=of N4] (N5) {结果分析与检验};
  \draw[arrowMain] (N1) -- (N2);
  \draw[arrowMain] (N2) -- (N3);
  \draw[arrowMain] (N3) -- (N4);
  \draw[arrowMain] (N4) -- (N5);
\end{tikzpicture}
\end{center}

\newpage

% ---- 场景 2：判断分支 ----
\sceneTitle{场景 2：判断分支流程图}
\begin{center}
\begin{tikzpicture}[node distance=\nodeDist and \colDist]
  \node[startstop] (S) {开始};
  \node[process, below=of S] (P1) {步骤一：初始数据计算};
  \draw[arrowMain] (S) -- (P1);
  \node[decision, below=of P1, yshift=-0.3cm] (D1) {条件判断\\（如收敛？精度满足？）};
  \draw[arrowMain] (P1) -- (D1);
  \node[process, right=of D1, xshift=1.0cm] (P2) {调整参数/迭代更新};
  \node[process, below=of D1, yshift=-0.3cm] (P3) {输出中间结果};
  \draw[arrowMain] (D1) -- node[edgeLabelYes, right] {否} (P2);
  \draw[arrowMain] (P2.east) -- ++(0.8,0) |- (P1.east);
  \draw[arrowMain] (D1) -- node[edgeLabelNo, left]  {是} (P3);
  \node[process, below=of P3] (P4) {后续处理};
  \draw[arrowMain] (P3) -- (P4);
  \node[startstop, below=of P4] (E) {结束};
  \draw[arrowMain] (P4) -- (E);
\end{tikzpicture}
\end{center}

\newpage

% ---- 场景 3：技术路线图 ----
\sceneTitle{场景 3：技术路线图}
\begin{center}
\begin{tikzpicture}[node distance=\nodeDistWide and \colDist]
  \node[startstop, minimum width=6cm, font=\song\small\bfseries] (T) {总体目标\\建立药材烘干过程的传热传质耦合模型};
  \node[process, below left=of T, xshift=-2.2cm, minimum width=2.8cm] (Q1) {问题一\\单颗粒干燥动力学};
  \node[process, below left=of T, xshift=0cm,   minimum width=2.8cm] (Q2) {问题二\\固定床堆积传热};
  \node[process, below right=of T, xshift=0cm,  minimum width=2.8cm] (Q3) {问题三\\连续带式干燥机};
  \node[process, below right=of T, xshift=2.2cm,minimum width=2.8cm] (Q4) {问题四\\分段温控优化};
  \draw[arrowMain] (T.south) -- ++(0,-0.3) -| (Q1.north);
  \draw[arrowMain] (T.south) -- ++(0,-0.3) -| (Q2.north);
  \draw[arrowMain] (T.south) -- ++(0,-0.3) -| (Q3.north);
  \draw[arrowMain] (T.south) -- ++(0,-0.3) -| (Q4.north);
  \node[subprocess, below=of Q1] (M1) {菲克扩散定律\\Page 薄层模型};
  \node[subprocess, below=of Q2] (M2) {多孔介质传热\\Ergun 压降方程};
  \node[subprocess, below=of Q3] (M3) {一维稳态模型\\热质耦合平衡};
  \node[subprocess, below=of Q4] (M4) {多目标优化\\遗传算法 / PSO};
  \draw[arrowSub] (Q1) -- (M1); \draw[arrowSub] (Q2) -- (M2);
  \draw[arrowSub] (Q3) -- (M3); \draw[arrowSub] (Q4) -- (M4);
  \node[subprocess, below=of M1, fill=white, draw=gray] (R1) {干燥速率曲线\\水分比变化规律};
  \node[subprocess, below=of M2, fill=white, draw=gray] (R2) {温度分布\\风速影响规律};
  \node[subprocess, below=of M3, fill=white, draw=gray] (R3) {出料含水率\\干燥效率};
  \node[subprocess, below=of M4, fill=white, draw=gray] (R4) {最优温控策略\\能效评估};
  \draw[arrowSub] (M1) -- (R1); \draw[arrowSub] (M2) -- (R2);
  \draw[arrowSub] (M3) -- (R3); \draw[arrowSub] (M4) -- (R4);
  \node[startstop, below=of R2, xshift=2cm, yshift=-0.5cm, minimum width=8cm, font=\song\small\bfseries] (Final) {综合最优烘干工艺方案};
  \draw[arrowMain] (R1.south) -- ++(0,-0.2) -| (Final.north);
  \draw[arrowMain] (R2.south) -- ++(0,-0.2) -| (Final.north);
  \draw[arrowMain] (R3.south) -- ++(0,-0.2) -| (Final.north);
  \draw[arrowMain] (R4.south) -- ++(0,-0.2) -| (Final.north);
\end{tikzpicture}
\end{center}

\newpage

% ---- 场景 4：并列对比 ----
\sceneTitle{场景 4：并列对比流程图}
\begin{center}
\begin{tikzpicture}[node distance=\nodeDist and 2.5cm]
  \node[startstop, minimum width=7cm] (S) {同一输入数据};
  \node[process, below left=of S, xshift=-1.2cm, minimum width=3.0cm] (A1) {方案 A\\（如：有限差分法）};
  \node[process, below right=of S, xshift=1.2cm, minimum width=3.0cm] (B1) {方案 B\\（如：有限元法）};
  \draw[arrowMain] (S.south) -- ++(0,-0.3) -| (A1.north);
  \draw[arrowMain] (S.south) -- ++(0,-0.3) -| (B1.north);
  \node[subprocess, below=of A1] (A2) {求解过程描述};
  \node[subprocess, below=of B1] (B2) {求解过程描述};
  \draw[arrowSub] (A1) -- (A2); \draw[arrowSub] (B1) -- (B2);
  \node[subprocess, below=of A2, fill=white, draw=gray] (A3) {结果 A：...};
  \node[subprocess, below=of B2, fill=white, draw=gray] (B3) {结果 B：...};
  \draw[arrowSub] (A2) -- (A3); \draw[arrowSub] (B2) -- (B3);
  \node[decision, below=of A2, xshift=2.5cm, yshift=-1.8cm, minimum width=4.0cm] (D) {对比分析\\哪个更优？};
  \draw[arrowMain] (A3.south) -- ++(0,-0.3) -| (D.west);
  \draw[arrowMain] (B3.south) -- ++(0,-0.3) -| (D.east);
  \node[startstop, below=of D, yshift=-0.3cm, minimum width=4.5cm, font=\song\small\bfseries] (E) {结论：推荐方案 X};
  \draw[arrowMain] (D) -- (E);
\end{tikzpicture}
\end{center}

\newpage

% ---- 场景 5：多阶段递进 ----
\sceneTitle{场景 5：多阶段递进流程图}
\begin{center}
\begin{tikzpicture}[node distance=\nodeDist and \colDist]
  \node[startstop] (S) {开始};
  \node[process, below=of S] (P1) {阶段一：全局预搜索};
  \draw[arrowMain] (S) -- (P1);
  \node[subprocess, below=of P1] (P1sub) {用粗网格/大步长扫描参数空间};
  \draw[arrowSub] (P1) -- (P1sub);
  \node[process, below=of P1sub, yshift=-0.2cm] (P2) {阶段二：局部精细优化};
  \draw[arrowMain] (P1sub) -- (P2);
  \node[subprocess, below=of P2] (P2sub) {在粗最优附近用梯度下降精细搜索};
  \draw[arrowSub] (P2) -- (P2sub);
  \node[decision, below=of P2sub, yshift=-0.3cm] (D) {精度满足？};
  \draw[arrowMain] (P2sub) -- (D);
  \node[subprocess, right=of P2sub, xshift=1.2cm] (adj) {缩小步长 / 扩大范围};
  \draw[arrowMain] (D) -- node[edgeLabelYes, above] {否} (adj);
  \draw[arrowSub] (adj.north) -- ++(0,1.0) -| (P2.east);
  \node[process, below=of D, yshift=-0.3cm] (P3) {阶段三：敏感性分析};
  \draw[arrowMain] (D) -- node[edgeLabelNo, left] {是} (P3);
  \node[startstop, below=of P3] (E) {输出最优解 + 置信区间};
  \draw[arrowMain] (P3) -- (E);
\end{tikzpicture}
\end{center}

\newpage

% ---- 场景 6：数据流图 ----
\sceneTitle{场景 6：数据流图}
\begin{center}
\begin{tikzpicture}[node distance=\nodeDist and 1.8cm]
  \node[data] (D1) {附件1\\实验条件数据};
  \node[data, right=of D1] (D2) {附件2\\干燥过程记录};
  \node[data, right=of D2] (D3) {附件3\\物性参数};
  \node[process, below=of D2, yshift=-0.2cm, minimum width=3.5cm] (Proc) {数据清洗与融合\\（异常检测、插值、对齐）};
  \draw[arrowMain] (D1.south) -- ++(0,-0.2) -| (Proc.north);
  \draw[arrowMain] (D2.south) -- (Proc.north);
  \draw[arrowMain] (D3.south) -- ++(0,-0.2) -| (Proc.north);
  \node[subprocess, below=of Proc, fill=white, draw=gray] (Mid) {\song\footnotesize 标准化数据集\\（n=XXX, 字段=XX）};
  \draw[arrowSub] (Proc) -- (Mid);
  \node[process, below=of Mid, minimum width=3.5cm] (Model) {模型建立与求解};
  \draw[arrowMain] (Mid) -- (Model);
  \node[data, below left=of Model, xshift=-0.6cm] (Out1) {干燥速率曲线};
  \node[data, below right=of Model, xshift=0.6cm] (Out2) {最优工艺参数};
  \draw[arrowMain] (Model.south) -- ++(0,-0.3) -| (Out1.north);
  \draw[arrowMain] (Model.south) -- ++(0,-0.3) -| (Out2.north);
\end{tikzpicture}
\end{center}

\newpage

% ---- 场景 7：带数学公式 ----
\sceneTitle{场景 7：带数学公式标注的流程图}
\begin{center}
\begin{tikzpicture}[node distance=\nodeDistWide and \colDist]
  \node[startstop] (S) {开始：输入物料参数与操作条件};
  \node[process, below=of S, minimum width=6cm] (P1) {Step 1：干燥动力学方程 \\[3pt] $MR = \dfrac{M_t - M_e}{M_0 - M_e} = \exp(-k t^n)$};
  \draw[arrowMain] (S) -- (P1);
  \node[process, below=of P1, minimum width=6cm] (P2) {Step 2：传热方程 \\[3pt] $\rho c_p \dfrac{\partial T}{\partial t} = \nabla\cdot(k\nabla T) + \dot{q}$};
  \draw[arrowMain] (P1) -- (P2);
  \node[decision, below=of P2, yshift=-0.3cm, minimum width=5cm] (D) {残差 $< \varepsilon$ \\ 且迭代次数 $< N_{\max}$};
  \draw[arrowMain] (P2) -- (D);
  \node[process, right=of D, xshift=1.2cm] (Adj) {调整松弛因子};
  \draw[arrowMain] (D) -- node[edgeLabelYes, above] {否} (Adj);
  \draw[arrowSub] (Adj.north) |- (P2.east);
  \node[process, below=of D, yshift=-0.3cm, minimum width=5cm] (P3) {Step 3：后处理 \\[3pt] 输出 $T(x,t)$, $M(x,t)$, 干燥速率曲线};
  \draw[arrowMain] (D) -- node[edgeLabelNo, left] {是} (P3);
  \node[startstop, below=of P3] (E) {结束};
  \draw[arrowMain] (P3) -- (E);
\end{tikzpicture}
\end{center}

\newpage

% ---- 场景 8：分组/分模块 ----
\sceneTitle{场景 8：分组/分模块流程图}
\begin{center}
\begin{tikzpicture}[node distance=\nodeDist and 2.5cm]
  \node[groupbox, minimum width=5.2cm, minimum height=2.5cm] (GA) {};
  \node[annotation, above=0.2cm of GA.north, font=\song\footnotesize\bfseries] {数据层};
  \node[data, minimum width=3.5cm] at (GA) (DA1) {原始数据采集};
  \node[data, below=0.6cm of DA1, minimum width=3.5cm] at (GA) (DA2) {预处理与标准化};
  \draw[arrowSub] (DA1) -- (DA2);
  \node[groupbox, right=of GA, xshift=1.0cm, minimum width=5.2cm, minimum height=2.5cm] (GB) {};
  \node[annotation, above=0.2cm of GB.north, font=\song\footnotesize\bfseries] {模型层};
  \node[process, minimum width=3.5cm] at (GB) (MB1) {传热传质耦合模型};
  \node[subprocess, below=0.6cm of MB1, minimum width=3.5cm] at (GB) (MB2) {参数辨识与验证};
  \draw[arrowSub] (MB1) -- (MB2);
  \node[groupbox, right=of GB, xshift=1.0cm, minimum width=5.2cm, minimum height=2.5cm] (GC) {};
  \node[annotation, above=0.2cm of GC.north, font=\song\footnotesize\bfseries] {输出层};
  \node[subprocess, fill=white, draw=gray, minimum width=3.5cm] at (GC) (RC1) {可视化结果};
  \node[subprocess, fill=white, draw=gray, below=0.6cm of RC1, minimum width=3.5cm] at (GC) (RC2) {最优工艺方案};
  \draw[arrowSub] (RC1) -- (RC2);
  \draw[arrowMain] (DA2.east) -- (DA2.east -| GA.east) -- ++(0.6,0) |- (MB1.west);
  \draw[arrowMain] (MB2.east) -- (MB2.east -| GB.east) -- ++(0.6,0) |- (RC1.west);
\end{tikzpicture}
\end{center}

\end{document}